% Generated draft: requires mathematical and source review.
\documentclass[11pt]{article}
\usepackage[margin=25mm]{geometry}
\usepackage{fontspec}
\setmainfont{Latin Modern Roman}
\usepackage{amsmath,amssymb,mathtools,unicode-math}
\setmathfont{Latin Modern Math}
\usepackage{enumitem,needspace,xurl,hyperref}
\setlength{\parindent}{0pt}
\setlength{\parskip}{5pt}
\setlength{\emergencystretch}{3em}
\newcommand{\N}{\mathbb{N}}
\newcommand{\Z}{\mathbb{Z}}
\newcommand{\Q}{\mathbb{Q}}
\newcommand{\R}{\mathbb{R}}
\newcommand{\C}{\mathbb{C}}
\newcommand{\F}{\mathbb{F}}
\newcommand{\A}{\mathbb{A}}
\newcommand{\PP}{\mathbb{P}}
\newcommand{\E}{\mathbb{E}}
\newcommand{\eps}{\varepsilon}
\newcommand{\dd}{\,\mathrm{d}}
\begin{document}
\section*{UnsolvedMath ID: 801093}
\textbf{Source problem code:} catalogue entry
\subsection{Definitions and mathematical statement}
Let \(\operatorname{Aut}(\C^3)\) denote the group of polynomial maps \(f\colon\C^3\to\C^3\) having polynomial inverses. For such a map, the standard first dynamical degree is
\[
\lambda_1(f)=\lim_{n\to\infty}\bigl(\deg(f^n)\bigr)^{1/n},
\]
where \(\deg\) is total polynomial degree. The supplied excerpt uses \(\lambda(f)\) without defining it; the formulation below interprets it as \(\lambda_1(f)\), subject to confirmation from the source context.

A projective compactification of \(\C^3\) is a projective variety \(X\) together with a dense open embedding \(j\colon\C^3\hookrightarrow X\). Write \(D=X\setminus j(\C^3)\) for its boundary. The map \(j\circ f\circ j^{-1}\) determines an induced birational self-map \(\bar f\colon X\dashrightarrow X\).

\textbf{Source question, with unresolved local terminology.} For every \(f\in\operatorname{Aut}(\C^3)\) satisfying
\[
\lambda_1(f)^2>\lambda_1(f^{-1}),
\]
prove that there exist a projective compactification \(X\) of \(\C^3\) and a point \(p\) on its divisor at infinity such that \(\bar f\) admits a super-attracting fixed point at \(p\). In particular, the fixed-point assertion requires \(\bar f\) to be defined and holomorphic on a neighborhood of \(p\), with \(\bar f(p)=p\).

The supplied excerpt does not specify regularity assumptions on \(X\), the precise divisor structure of its boundary, or the meaning of ``super-attracting.'' At a smooth point \(p\), two possible local conditions are
\[
\mathrm d\bar f_p=0
\qquad\text{and}\qquad
\mathrm d\bar f_p\text{ is nilpotent}.
\]
These conditions are not equivalent in dimension three. The draft does not select between them or impose smoothness or normality on \(X\); the intended convention and admissible compactifications require source-context verification.

\textbf{Supplied motivation.} The excerpt reports that, under the same dynamical-degree inequality, an invariant valuation on \(\C[x_1,x_2,x_3]\) is known to exist by the reference labeled \([\mathrm{DF20}]\). It informally describes such a valuation as selecting an irreducible subvariety in every algebraic compactification of \(\C^3\), compatibly across compactifications. The excerpt supplies neither the precise invariance convention nor the compatibility statement. It also says that extending an earlier construction to a larger class of polynomial automorphisms of \(\C^3\) is unclear; that earlier construction is not included in the record.
\subsection{Short English statement}
The supplied question asks whether every polynomial automorphism of \(\C^3\) satisfying \(\lambda(f)^2>\lambda(f^{-1})\) has a projective compactification on which its induced birational map has a super-attracting fixed point at infinity. The dynamical-degree notation, admissible compactifications, and exact local meaning of super-attraction require confirmation. The registry labels the problem open; this draft does not independently verify that status.
\subsection{Sources}
\begin{itemize}
\item Local UnsolvedMath registry, permanent ID \(801093\), exact source problem code \texttt{catalogue entry}, title ``Henon Maps: Open Problems --- Question 45.'' The record attributes the problem to Berger and Bedford, with proposed year 2023, and supplies the status \texttt{open}.
\item Berger and Bedford, \emph{H\'enon Maps: A List of Open Problems} (2023), Question 45. Supplied source URL: \url{https://arxiv.org/abs/2312.03907}. The record reports access on 2026-10-16 and extraction by \texttt{arxiv\_latex}; these are supplied provenance statements, not independently verified access or literature assessments.
\item Reference \([\mathrm{DF20}]\), cited in the supplied excerpt for existence of an invariant valuation under \(\lambda(f)^2>\lambda(f^{-1})\). No authors, title, bibliographic details, or separate link for this reference are supplied, so its identity and precise theorem remain unverified.
\item Supplied upstream revision: \texttt{d371ee917ac9647fbb06cd6fc529deec6f3dda82}.
\end{itemize}
\end{document}
